
\documentclass[12pt]{article}

\usepackage{caption}
\usepackage{float}
\usepackage{hyperref}
\usepackage{acronym}

\author{Samuel Brunner}
\date{\today}
\title{Technische Dokumentation Kühlschrankmanager KM3003}

\begin{document}

\maketitle

\newpage
\tableofcontents

\newpage

\section{Einleitung}
Der Kühlschrankmanager im Vertretungszimmer der Fachschaft ersetzt die Strichliste, welche früher am Kühlschrank hing, und die Getränkekäufe der Vereinsmitglieder erfasste.

Früher wurden auf der Liste für jedes erworbene Getränk Striche in der Zeile des jeweiligen Mitglieds und in der Spalte des Getränks eingetragen. Regelmäßig zählte man die Striche, ermittelte so die Saldos der Mitglieder und erstellte eine neue Liste.

Der KM3003 ist die Weiterentwicklung der bisherigen Kühlschrankmanager (KM3000 bis KM3002). Der KM3002 basierte auf einem Raspberry Pi mit Touchscreen und NFC-Reader in einem 3D-gedruckten Gehäuse. Mitglieder wurden per Studentenausweis mit NFC identifiziert, und Käufe über den Touchscreen ausgewählt. Die Daten wurden in einer MySQL-Datenbank gespeichert, die jedoch nur negative Salden ermöglichte und so exakte Abrechnungen erzwang. Zudem gab es Verbesserungsbedarf in Codequalität, Zuverlässigkeit und Geschwindigkeit.

Der KM3003 erlaubt nun die Identifizierung von Mitgliedern und Produkten über das Scannen von Barcodes. Mitglieder werden über die Barcodes auf den Studentenausweisen erkannt, Produkte über EAN-Codes. Außerdem erlaubt er nun auch positive Saldos, sodass Mitglieder flexibel Guthaben einzahlen können.

Zusammengefasst ist der KM3003 ein Self-Checkout Point-of-Sale-Gerät mit Touchscreen und Barcode-Scanner.

\section{Plattform und Systematchitektur}
Als Basis für den Kühlschrankmanager wurde ein Microsoft Surface 3 mit dem zugehörigen Dock verwendet.
Das Surface 3 ist ein Tablet-Computer mit einem zuverlässigen kapazitativen Multi-Touchscreen und solider Treiberunterstützung unter Windows.
Als Betriebssystem wird ein unverändertes Windows 11 Home verwendet, um eine gut getestete und stabile Basis zu bieten. Auch die Verfügbarkeit von Updates in der Zukunft wird so maximiert.

Weiterhin bieten die meisten Microsoft Surface Geräte einen Kioskmodus der spezielle Lösungen für Geräte mit festem und orstunveränderlichem Verwendungszweck bietet.
So optimiert der Kioskmodus das Ladeverhalten für Geräte, die permanent mit Strom versorgt werden, um die Akkulebensdauer zu verlängern.
Weiterhin kann der Zugang zu Anwendungen, die nicht dem gedachten Verwendungszweck entsprechen, beschränkt werden.

Zum scannen wird ein Freihandscanner (Modell DS9208) der Firma Zebra verwendet. Dieser wird per USB angeschlossen, meldet sich aber unter Windows, als ein über serielle Schnittstelle angebundenes Gerät.
Dies erfolgt unter Verwendung des \ac{USB CDC} Treibers.

Das Surface speichert keinerlei Daten. Hierfür wird eine externe MySQL-Datenbank verwendet, welche Nutzer, Produkte, Ein- bzw. Auszahlungen und Einkäufe enthält.

\section{Software}
Als Implementierungssprache wurde objektorientiertes Python3 verwendet. Python ist weit verbreitet und einfach zu lernen, was Wartungsarbeiten oder die Erweiterung mit neuen Features für zukünftige Mitglieder erleichtert. 

Weiterhin können Python Anwendungen auf Linux und Windows Platformen ausgeführt werden, sodass bei Bedarf auch das bisher verwendete Raspberry Pi mit externem Touchscreen weiter verwendet werden könnte. 

Auch die Verfügbarkeit von bestehenden Bibliotheken für \ac{GUI}, Datenbankanbindung und serielle Schnittstellen ist optimal.

\subsection{Verwendete Bibliotheken und Frameworks}

% - Lizenzmodell
%     PySimpleGUI; PySimpleGUI license
%         * "Hobbyist Developer" means any individual who uses PySimpleGUI for
%         development purposes solely for either or both of the following: (1) personal 
%         (e.g., not on behalf of an employer or other third party), Non-Commercial
%         purposes; or (2) Non-Commercial educational or learning purposes (1 and 2
%         together, the "Permitted No-cost Purposes").
%         Lizenz ist aktuell so und erlaubt uns die Nutzung, weil wir keinen Gewinn machen. Ziel der Kasse ist 0. Bleibt das so?
%         Die Andere Option wäre \$99 zu zahlen. (Jährlich um auch die zukünftigen Releases zu nutzen). Also eigtl \$99 pro Jahr
%     tkinter: Python Software Foundation License
%         * The authors hereby grant permission to use, copy, modify, distribute, and license this software and its documentation for any purpose, provided that existing copyright notices are retained in all copies and that this notice is included verbatim in any distributions. 
%         Python Software Foundation License aggreement: Alows unrestricted distribution 
%      qt: GPL license auch ok für uns... muss dann nur auch als GPL lizensiert werden
% - Dokumentation:
%     PySimpleGUI: 
%         extrem auf den Hobbyist ausgelegt, sehr beispiellastig
%         meiner meinung nach furchtbar zu lesen
%         hätte lieber eine ordentliche API spezifikation, als 10000 beispielprogramme
%     tkinter: standard doku
%         weniger an den Hobbyisten gerichtet; 
%         für den geübten aber einfacher das zu finden was man sucht (Werkzeuge, wie configure() )
%     qt: schlechte Doku und 

% - Architektur und Objekte:
%     tkinter: 
%         everything is a widget
%         widgets are organized in a hierarchy
%         widgets have configuration options
%         the placement options is managed by a geometry manager 
%         the GUI reacts to user inputs and changes from my program and refreshes the interface
%     pysimplegui:
%         everything is an element, which is also just a pyhton object
%         layout is done by filling rows with elements and pushing them vertically to where you want them
%         some elements can contain other elements, which are called container elements
%         elements have configuration options
%     qt:
%         offers far to much customizability for out purposes


% conclusion: If I could do it again, I would use tkinter, because of the license that im sure of that it wont change in the future, the better documentation (Pysimplegui docs is annoying af) and the more professional appearance

Bei der Auswahl eines GUI-Frameworks für Python-Projekte bieten PyQt, Tkinter und PySimpleGUI unterschiedliche Stärken und Einschränkungen in Bezug auf Lizenzmodell, Dokumentation und Architektur. 

Tkinter und PyQT bieten Python Schnittstellen zu den externen GUI toolkits Tk und Qt. PySimpleGUI hingegen baut auf solchen toolkits auf und legt eine weitere Abstraktionsebene darüber um das Erstellen von GUIs über diese Toolkits hinweg zu vereinheitlichen und weiter zu vereinfachen.

PySimpleGUI unterliegt seit Version 5 einer projekteigenen Lizenz, welche die kostenlose Nutzung für Hobbyentwickler und Bildungszwecke erlaubt, sofern keine kommerzielle Nutzung stattfindet. Tkinter und PyQT sind GPL kompatibel bzw. direkt mit GPL linenziert. Bei Tkinter und PyQt ist nicht davon auszugehen, dass sich an der Lizenzierung etwas ändert. Bei PySimpleGUI hingengen muss jedes Jahr wieder eine neue Hobbyisten- Lizenz über einen Account gelöst und wieder im Projekt eingepflegt werden. In Anbetracht der bisherigen Änderungen und dem zusätzlichen Arbeitsaufwand die Lizenz jährlich zu aktualisieren ist PySimpleGUI lizenztechnisch die schlechteste Wahl.

Die Dokumentation der drei Frameworks variiert stark in der Zielgruppe und Benutzerfreundlichkeit. PySimpleGUI ist stark auf Hobbyentwickler ausgelegt und verwendet viele Beispielprogramme, um Funktionen zu erklären. Dies kann für Anfänger hilfreich sein, erweist sich jedoch für erfahrene Entwickler, welche eher eine formale API-Spezifikation bevorzugen, als mühsam und unstrukturiert. Tkinter und PyQT bieten eine standardisierte Dokumentation, die weniger auf Anfänger zugeschnitten ist, aber erfahrenen Entwicklern eine effizientes Nachschlagewerk bietet.

Architektonisch ähneln sich die Frameworks in der Handhabung und Flexibilität ihrer Komponenten. Tkinter bezeichnet GUI-Komponenten als Widget, PySimpleGUI als Elements und PyQt

das in einer Hierarchie organisiert ist. Widgets lassen sich über umfangreiche Konfigurationsoptionen anpassen und über einen Geometriemanager präzise platzieren, was eine flexible Gestaltung und Reaktionsfähigkeit der GUI ermöglicht. PySimpleGUI hingegen strukturiert die GUI über sogenannte Elemente, die als einfache Python-Objekte modelliert sind. Der Aufbau der Oberfläche erfolgt, indem Elemente in Zeilen gefüllt und vertikal gestapelt werden. Container-Elemente ermöglichen eine einfache Gruppierung, jedoch ist die Flexibilität der Gestaltung im Vergleich zu Tkinter eingeschränkt, was bei komplexen Layoutanforderungen zu Einschränkungen führen kann. PyQt bietet ein sehr umfangreiches Spektrum an Anpassungsmöglichkeiten, das weit über die Anforderungen einfacher Anwendungen hinausgeht. Dies kann für spezialisierte und detaillierte Benutzeroberflächen vorteilhaft sein, jedoch auch die Entwicklung durch den hohen Grad an Komplexität erschweren.


Für einfache Projekte und eine schnelle Umsetzung bietet PySimpleGUI durch seine einfache, auf Hobbyentwickler ausgerichtete Architektur eine zugängliche Option. Bei langfristigen Projekten oder Anwendungen, die stabil und flexibel gestaltet sein sollen, eignet sich Tkinter aufgrund seiner stabilen Lizenzierung, der gut strukturierten Dokumentation und der flexiblen Architektur besser. PyQt hingegen ist ideal für Projekte, die sehr anpassbare und spezialisierte Benutzeroberflächen erfordern und keine Probleme mit der GPL-Lizenz haben.



mysql connector

this is a test


\subsection{Datenbankstruktur}
TODO an Normalformen angelehnt
atomare Felder

\subsection{Testing}
Unit tests schreinben mit pyunit test Framework





\section{Deployment}
- deployment:
    surface, ansbile for km3003 standalone
    cyclical restart windows task scheduler
    mysql db

\section{Herausforderung bei der Implementierung}
    - wiedreraufbau von Verbindunsabbrüchen (Datenbank und Scanner) zur Laufzeit

\section{Ausblick und Verbesserungsvorschläge}
    UI Framework is kacke
    Sounds

\section{Inline documetation anhängen}
\appendix

\begin{acronym}[OTH R]
    \acro{USB CDC}{USB communications device class}
    \acro{GUI}{graphical user interface}
\end{acronym}

\clearpage

\listoffigures

\clearpage

\begin{thebibliography}{99}
\bibitem{bl}Microsoft Kiosk Mode \url{https://learn.microsoft.com/en-us/windows/iot/iot-enterprise/customize/kiosk-mode}
\bibitem{bl}What is infrastructure as code (IaC)? \url{https://learn.microsoft.com/en-us/devops/deliver/what-is-infrastructure-as-code}
\bibitem{bl}tkinter — Python interface to Tcl/Tk \url{https://docs.python.org/3/library/tkinter.html}
\bibitem{bl}PyQt vs. Tkinter — Which Should You Choose for Your Next GUI Project? \url{https://www.pythonguis.com/faq/pyqt-vs-tkinter/}
\bibitem{bl}PyQT6 Reference Guide \url{https://www.riverbankcomputing.com/static/Docs/PyQt6/index.html}



\end{thebibliography}

\clearpage
\end{document}